\setlength{\headheight}{14.5pt}

% --- 1. ENCODAGE ET LANGUE (À mettre en premier) ---
\usepackage[T1]{fontenc}
\usepackage[utf8]{inputenc}
\usepackage[french]{babel}

% --- 2. POLICES ET SYMBOLES ---
% fourier : police élégante pour les maths (lmodern supprimé car doublon)
\usepackage{fourier}
\usepackage{amsfonts, amssymb, amsthm}
\usepackage[right]{eurosym}
\usepackage[np]{numprint}
\usepackage[nointegrals]{wasysym}

% --- 3. GÉOMÉTRIE ET MISE EN PAGE ---
\usepackage[left=1cm,right=1cm,top=1.5cm,bottom=1.5cm]{geometry}
\usepackage[dvipsnames,svgnames]{xcolor}
\usepackage[most,skins,breakable]{tcolorbox}
\usepackage{tagging}

% --- 4. GRAPHISMES ET DIAGRAMMES ---
\usepackage{graphicx}
\usepackage{pgfplots}
\pgfplotsset{compat=1.18}
\usetikzlibrary{babel, shapes, fit} % 'babel' évite les bugs avec TikZ

% --- 5. TABLEAUX ---
\usepackage{array, booktabs, multirow, tabularx, longtable}

% --- 6. STRUCTURE ET TITRES ---
% titlesec est plus moderne que sectsty, on garde uniquement celui-ci
\usepackage{titlesec}
\usepackage{enumitem}
\usepackage[normalem]{ulem} % 'normalem' évite que \emph devienne souligné

% --- 7. HYPERLIENS (toujours en dernier ou presque) ---
\usepackage{hyperref}

\setlength{\parindent}{0pt}
\frenchbsetup{StandardItemLabels=true}

% Le package ultime
\usepackage{ProfCollege}


% =================================================================
% COMMANDES GÉNÉRALES
% =================================================================

% Une ligne de points (usage : \points[15])
\newcommand{\points}[1][1]{%
    \foreach \n in {1,...,#1} { \ldots}%
}

% Plusieurs lignes de points (usage : \lignesPoints{3})
\newcommand{\lignesPoints}[1]{%
    \foreach \n in {1,...,#1} {%
        \dotfill \newline
    }%
}


% =================================================================
% ZONES RÉPONSE
% =================================================================

% Zone réponse universelle (usage : \begin{ZoneReponse}{4cm})
\newtcolorbox{ZoneReponse}[1]{
    enhanced,
    colback=white,
    colframe=gray!50,
    arc=2mm,
    boxrule=0.5mm,
    width=\textwidth,
    height=#1,
    valign=top,
    fonttitle=\bfseries,
    colbacktitle=gray!20,
    coltitle=black,
    attach title to upper,
    after skip=10pt plus 2pt minus 2pt
}


% =================================================================
% COMMANDES MATHÉMATIQUES COLLÈGE
% =================================================================

% Ensembles de nombres
\newcommand{\N}{\ensuremath{\mathbf{N}}}
\newcommand{\Z}{\ensuremath{\mathbf{Z}}}
\newcommand{\D}{\ensuremath{\mathbf{D}}}
\newcommand{\Q}{\ensuremath{\mathbf{Q}}}
\newcommand{\R}{\ensuremath{\mathbf{R}}}

% Valeur absolue
\newcommand{\abs}[1]{\left\lvert #1 \right\rvert}

% Arc géométrique (usage : $\arc{AB}$)
% Note : utilise la fonte tipa pour le symbole arc
\makeatletter
\DeclareFontFamily{U}{tipa}{}
\DeclareFontShape{U}{tipa}{m}{n}{<->tipa10}{}
\newcommand{\arc@char}{{\usefont{U}{tipa}{m}{n}\symbol{62}}}
\newcommand{\arc}[1]{\mathpalette\arc@arc{#1}}
\newcommand{\arc@arc}[2]{%
  \sbox0{$\m@th#1#2$}%
  \vbox{\hbox{\resizebox{\wd0}{\height}{\arc@char}}\nointerlineskip\box0}%
}
\makeatother

% Comparaison : boîte vide entre deux nombres pour écrire le symbole
% Usage : \compare{102}{120}  ou  \compare[0.5]{\np{132999999}}{\np{133000000}}
% Argument optionnel : ignoré (compatibilité ascendante), la boîte s'adapte
\newcommand{\compare}[3][0]{%
    #2 \kern0.4em
    \tikz[baseline=-0.6ex]
        \node[rectangle, draw=gray!60, thick, dotted, minimum height=0.6cm,
              minimum width=0.6cm, inner sep=2pt] {};%
    \kern0.4em #3%
}

% Encadrement : deux boîtes vides autour d'un nombre
% Usage : \encadre{\np{1389}}  ou  \encadre[1.8]{\np{652176}}
% Argument optionnel : largeur des boîtes en cm (défaut 1.5)
\newcommand{\encadre}[2][1.5]{%
    \tikz[baseline=-0.6ex]
        \node[rectangle, draw=gray!60, thick, dotted, minimum height=0.6cm,
              minimum width=#1cm, inner sep=2pt] {};%
    \kern0.3em $<$ #2 $<$ \kern0.3em
    \tikz[baseline=-0.6ex]
        \node[rectangle, draw=gray!60, thick, dotted, minimum height=0.6cm,
              minimum width=#1cm, inner sep=2pt] {};%
}

% Diagramme de proportionnalité simple
% Usage : \diagrammesimple{A}{B}{×3}
\newcommand{\diagrammesimple}[3]{%
\begin{tikzpicture}[baseline=(A.base)]
    \node(A) at (0,0) {#1};
    \node(B) at (5,0) {#2};
    \draw[->, dashed, line width=0.8mm, green!50!black]
        (A) to[bend left=20] node[midway, above=2pt] {#3} (B);
\end{tikzpicture}}

% Diagramme avec flèche aller-retour
% Usage : \diagrammesimplereciproque{A}{B}{×3}{÷3}
\newcommand{\diagrammesimplereciproque}[4]{%
\begin{tikzpicture}[baseline=(A.base)]
    \node(A) at (0,0) {#1};
    \node(B) at (5,0) {#2};
    \draw[->, dashed, line width=0.8mm, green!50!black]
        (A) to[bend left=20] node[midway, above=2pt] {#3} (B);
    \draw[->, dashed, line width=0.8mm, blue!50!black]
        (B) to[bend left=20] node[midway, below=2pt] {#4} (A);
\end{tikzpicture}}